% S23B, DECISIONS item 150. Section 5 was a stub through S18 to S23 under item 131.
% The body below is the text supplied with the S23B prompt, converted verbatim.
\section{Kill conditions}\label{sec:conditions}

% The HAR family the attenuation condition operates on is \citet{corsi2009}
% and its quarticity-augmented form \citep{bpq2016}.
\nocite{corsi2009,bpq2016}

Thirteen conditions were pre-registered, each with its threshold fixed before the
measurement it adjudicates and its null abstract drafted before any result was
seen. Table~\ref{tab:kill} gives all thirteen. A condition fires when the null
holds, which is to say when the measured quantity falls inside the threshold and
reliability does not change the decision.

Two of the thirteen concern the measurement itself. K2 asks whether any
reliability estimator is grid-invariant and fires, in the sense that none is: the
ratio of largest to smallest implied $\operatorname{Var}(\log IV)$ across a grid
runs from $1.05$ for the best estimator to $1.97$ for the worst. K3 asks whether % k_table.csv: K2 margin
the scaling result survives its own uncertainty and stands, with the reference
inside the interval in $0$ of $54$ fits. The remaining eleven are applications. % numbers.csv: n_proxy_ref_inside_95, n_proxy_fits

\begin{table}[!ht]\centering\footnotesize
\setlength{\tabcolsep}{4pt}\renewcommand{\arraystretch}{1.15}
\setlength{\abovetopsep}{6pt}
\caption{Kill conditions and their determinations.\protect\footnotemark{} Generated from the determination artifacts listed in \texttt{paper/k\_table.csv}; the margin column gives the measured quantity against its pre-registered threshold, and the final column records whether the application was selected before or after the first-order criterion was articulated.}\label{tab:kill}
\begin{tabular}{@{}ll>{\raggedright\arraybackslash}p{3.3cm}>{\raggedright\arraybackslash}p{3.0cm}>{\raggedright\arraybackslash}p{4.3cm}c@{}}
\toprule
Paper & Session & Content & Determination & Margin & Selected \\
\midrule
K1 & K1/K2 & reliability correction does not change MCS composition & INDETER\-MINATE & excess 20.8pp on clean geometry, not tracking lambda & before \\  % source: artifacts/s09-application/phase12_summary.json [K2, excess_rth]
K2 & K2 (grid) & no reliability estimator is grid-invariant & FIRES & lambda*Var(log RV) max/min ratio 1.05 best, 1.97 worst & before \\  % source: docs/specs/SPEC-obs-space-vol-eval.md [section 7 null abstract]
K3 & K3 (scaling) & proxy-error scaling inconsistent with sampling theory & STANDS & reference inside the 95\% interval in 0 of 54 proxy fits & before \\  % source: artifacts/s11-extensions/phase7_summary.json [ref_inside_by_proxy]
K4 & K3 (sizing) & sizing consequence null & FIRES & max R2-vs-R3 relative TE difference 1.008\% vs 5\% & before \\  % source: artifacts/s09-application/phase7_k3.json [extended.K3]
K5 & K4 & risk-limit breaches & stop-out FIRES; cap UN\-TESTED & spurious 0.966\% vs 1\%; cap bound 0 of 2524 & before \\  % source: artifacts/s12-correction/phase4_k4.json [stop_out, leverage_cap]
K6 & K5 & inverse-MSE combination weights & FIRES & max relative TE difference 0.50\% vs 5\% & before \\  % source: artifacts/s11-extensions/phase8910_summary.json [K5]
K7 & K6 & variance-to-volatility convexity adjustment & DOES NOT FIRE & overstatement 5.9-134.3\% vs 5\% & before \\  % source: artifacts/s12-correction/phase1_k6.json [K6]
K8 & K7 & inverse-volatility risk parity & FIRES & mean $|$dw$|$ 0.000145 vs 0.02 & before \\  % source: artifacts/s13-extension/phase2_k7.json [K7]
K9 & K8 & regime misclassification at a median threshold & DOES NOT FIRE & analytic 7.9-13.6\% vs 5\% & after \\  % source: artifacts/s14-applications/phase1_k8.json [K8]
K10 & K9 & HAR persistence attenuation & INDETER\-MINATE & daily shift 4.8-116\% vs 10\%, spanning the threshold on the Sigma\_E choice & after \\  % source: artifacts/s15-confounds/s15_summary.json [K9]
K11 & K10 & Hurst exponent bias from proxy noise & DOES NOT FIRE & shift 0.026-0.213 vs 0.02 on a fixed lag window & after \\  % source: artifacts/s15-confounds/s15_summary.json [K10]
K12 & K11 & regime classification under an observable-side correction & DOES NOT FIRE & 0.00pp from the correction; effect is A3 not being A2 & after \\  % source: artifacts/s16-regime/phase4_k11_final.json [K11]
K13 & K12 & measurement error in the emission & FIRES & 0.00pp maximum reduction at every scaling, floor binding in 63\% of windows & after \\  % source: artifacts/s17-emission/phase345_summary.json [K12]
\bottomrule
\end{tabular}
\end{table}
\footnotetext{The session record uses overlapping labels. Item~61 relabelled the MCS-composition condition K2 while the specification already used K2 for grid invariance, and two distinct conditions were both called K3. The paper renumbers K1--K13 in the order shown; the repository retains the original labels, given here in the second column.}


\subsection{What separates them}

Proxy noise is second-order where the loss surface is smooth in the corrected
quantity, or where the contamination averages away as the estimation window
lengthens. It is first-order where the quantity of interest depends on the
variance of the estimate rather than its level, and where more data therefore
does not reduce the contamination. The distinction is not about the size of
$\lambda$, which is the same in both cases, and it predicts which side a new
application falls on.

\subsection{The second-order cases}

Volatility-target sizing is quadratic in the shrinkage weight near its optimum, so
a first-order error in $\lambda$ produces a second-order loss: replacing the
textbook weight with the measured one moves out-of-sample tracking error by at
most $1.008$ percent against a $5$ percent threshold, in the same direction in % k_table.csv: K4 margin
every cell, at every point of an eightfold cost sweep. Inverse-MSE combination
weights move by $0.50$ percent in the same metric, because the proxy noise % k_table.csv: K6 margin
variance is common to every model in the set and cancels from the relative
weights to first order. Inverse-volatility risk parity moves the allocation by a
mean absolute weight deviation of $0.000145$ against a $0.02$ threshold, and the % k_table.csv: K8 margin
reason generalises: the bias in $\mathbb{E}[1/\hat{\sigma}]$ scales as one over
the averaging window, so monthly rebalancing on a twenty-one-session estimate
divides it by twenty-one. Risk limits fire on the stop-out at a spurious breach
rate of $0.966$ percent against $1$ percent, and we report the leverage cap as % k_table.csv: K5 margin
untested rather than as passing, since it bound in $0$ of $2{,}524$ decision
points and would first bind at $1.434$ times a ten percent target. % numbers.csv: k5_cap_first_binds

\subsection{The first-order cases}

The variance-to-volatility convexity adjustment depends on
$\operatorname{Var}(\log IV)$ directly, and substituting the proxy overstates it
by between $5.9$ and $134.3$ percent, which exceeds its $5$ percent threshold in % numbers.csv: convex_overstatement[NQ/GLOBEX/1day], convex_overstatement[ES/RTH/30min]
every cell. Averaging over more windows improves the estimate of a contaminated
population quantity and does not decontaminate it. A regime threshold placed at
the median has no curvature to absorb the error, and misclassification runs
between $7.9$ and $13.6$ percent against a $5$ percent threshold. Proxy noise in % numbers.csv: mis_analytic[NQ/GLOBEX/insample], mis_analytic[ES/GLOBEX/insample]
the estimation of the Hurst exponent enters as a nugget that does not vanish at
short lags, and correcting for the measured $\lambda$ shifts $H$ by between
$0.026$ and $0.213$ on a fixed lag window against a $0.02$ threshold. % k_table.csv: K11 margin

\subsection{The two indeterminate conditions}

K1 asks whether the reliability correction changes model confidence set
composition. Composition differs in $20.8$ percent of comparisons on the clean % numbers.csv: k1_excess_rth
geometry after subtracting a placebo scheme that conditions on a variable with no
proxy content, so the placebo accounts for roughly seventy percent of the raw rate
and what survives it is the remaining twenty-one points. What remains % numbers.csv: k1_raw_rate_rth, k1_placebo_rate_rth, k1_excess_rth
does not track $\lambda$ across horizons, running $12.5$, $43.8$ and $6.3$ percent % numbers.csv: k1_excess[1day], k1_excess[1h], k1_excess[30min]
as reliability falls from $1.001$ to $0.856$ to $0.713$, which is not the ordering % numbers.csv: k1_lam_intercept[1day], k1_lam_intercept[1h], k1_lam_intercept[30min]
the mechanism predicts. We report it as indeterminate rather than as either
result. K10 asks whether the attenuation of HAR persistence coefficients is
material. The daily coefficient shifts by between $4.8$ and $116$ percent
depending on which construction supplies the measurement-error covariance, and % k_table.csv: K10 margin
the two constructions bracket the $10$ percent threshold from either side. The
bracket is the same one that leaves the composition of $A M^{b}$ unresolved in
Section 4.2, resurfacing inside an application.

\subsection{The criterion's record}

Eight conditions were selected before the criterion of Section 5.1 was
articulated and five of them fire. Five were selected after it and three return % numbers.csv: n_conditions_before, n_fire_before, n_conditions_after, n_material_after
material effects. We note the ordering because the criterion was derived from the
first eight and tested on the last five, not because five of thirteen is a rate
worth reporting. Two of the five later conditions test whether the measured
coefficient can be inserted at all rather than whether it matters, and Section 6
takes those up separately.

